\chapter{Object-oriented, functional, procedural, and data-oriented design}

\section{Paradigms are lenses}
A programming paradigm emphasizes certain ways of representing state,
behaviour, and composition. Object-oriented design groups data with operations
and uses encapsulation and polymorphism. Functional design emphasizes expressions,
immutable values, and explicit transformation. Procedural design emphasizes
ordered commands and control flow. Data-oriented design arranges data for the
access patterns of the hardware and workload.

Real systems combine them. A functional core can sit inside an object-oriented
application; a procedural migration script can protect a data-oriented storage
layout. The question is whether the chosen representation makes the important
invariants and costs visible.

\section{Encapsulation and state}
Encapsulation means clients depend on a stable set of operations rather than
directly manipulating representation. It helps when the invariant is local and
the state has a meaningful lifecycle. It becomes harmful when a class is merely a
namespace for unrelated functions or when inheritance creates hidden coupling.

Prefer composition when behaviour varies independently. Inheritance can express
a genuine substitutability relation, but sharing code is not by itself a reason
to create a subtype. A payment gateway that supports `authorize` and `capture`
may be a port with several adapters; it need not form a deep class hierarchy.

\section{Functional core, imperative shell}
Pure functions are easier to test because their result depends on explicit input.
Side effects are still necessary for useful software. A practical split is:

\begin{itemize}
  \item imperative shell: read HTTP, clock, randomness, database, and network;
  \item functional core: validate, calculate totals, decide transitions, and
        construct commands from explicit values.
\end{itemize}

The Python \texttt{total\_cents} example follows this approach. It accepts data, applies
rules, and has no network or database dependency. The shell decides when and how
to call it. This does not eliminate integration tests; it makes the rule-level
tests fast and diagnostic.

\section{Procedures and data flow}
Procedural code can be the clearest form for a workflow. A migration or a shell
script is often easier to audit as a sequence of explicit steps than as a graph
of objects. The risk is accidental shared state and unclear error handling. Keep
preconditions, postconditions, and cleanup paths visible.

\section{Data-oriented thinking}
Performance depends on the whole path: algorithm, data layout, allocation,
cache behaviour, I/O, concurrency, and workload shape. Data-oriented design
starts with the hot operations and organizes data so they do not repeatedly
load irrelevant fields. It is useful in simulations, games, analytics, and
high-throughput infrastructure; it can make ordinary business code less
readable when applied without a measured bottleneck.

\section{Choosing a representation}
Use these questions:

\begin{enumerate}
  \item What state must be protected by an invariant?
  \item Which effects are unavoidable, and can they be isolated?
  \item Does behaviour vary by a stable capability or by a changing data value?
  \item Is the workload limited by CPU, memory, allocation, I/O, or coordination?
  \item What will a future reader need to change safely?
\end{enumerate}

\noindent\begin{tabularx}{\textwidth}{@{}p{0.24\textwidth}X@{}}
\toprule
Situation & Often useful starting point \\
\midrule
Local lifecycle invariant & Encapsulation around a small state machine \\
Many transformations and rules & Pure functions and immutable values \\
Explicit workflow or migration & Procedural sequence with checked steps \\
Measured hot path and predictable data & Data-oriented layout and profiling \\
\bottomrule
\end{tabularx}

\begin{exercise}
Represent an order cancellation in two ways: a mutable object and a pure
transition function. For each, state where invalid state is prevented and where
side effects occur. Which representation gives the clearer test failure?
\end{exercise}

\section{Chapter summary}
Paradigms are tools for making state, effects, and cost visible. Combine them
deliberately. Prefer the simplest representation that protects the important
invariants and fits the measured workload, and avoid turning stylistic taste into
an architectural law.
